% ch-discretization.tex -- numerical method: the grid, the ghost layers and the matrix rows
\chapter{離散化}
\label{ch:discretization}

本章では、Yuzenの格子について説明する。物理量を置ける位置、境界の外側の拡張のしかた、格子上の点と行列の行との対応を順に述べる。ここで導入する用語は\cref{ch:timeintegration,%
ch:convergence,ch:linearsolver,ch:reproducibility}で用いるので、物理項を追加する場合はまず本章を参照されたい。離散化の対象は\cref{ch:governing}の\cref{eq:height,eq:flux}である。格子の大きさを決めるキーは\cref{sec:ref-domain:grid}に示す。

\section{スタガード格子と八種類の配置型}
\label{sec:disc:stagger}

Yuzenは一様な矩形格子を用いる。格子線を$x = x_{\min} + i_x\,\Delta x$とすると、物理量を格納できる場所は各軸について二つしかない。

\begin{center}
\begin{tabular}{@{}l l@{}}
\code{0} & 格子線上（$x = x_{\min} + i_x\,\Delta x$） \\
\code{1} & 隣り合う格子線の中間（$x = x_{\min} + i_x\,\Delta x + \Delta x/2$） \\
\end{tabular}
\end{center}

軸が三つあり、それぞれに二通りの選択肢があるので、可能な位置は3ビットで表せ、全部で八通りになる。この3ビットを表すのが型\code{Stagger<SX,SY,SZ>}（\file{yuzen/src/grid/Stagger.h}）である。各位置の名前を\cref{tab:site-types}に示す。

\begin{yzntable}{八種類の配置型と3ビットによる符号化\label{tab:site-types}}{@{}l l L{0.46\textwidth}@{}}
\textsf{\bfseries ビット $(x,y,z)$} & \textsf{\bfseries 名前} &
\textsf{\bfseries 位置} \\ \midrule
0,0,0 & \code{Node}  & 格子点 \\
1,0,0 & \code{Xedge} & $x$方向の辺 \\
0,1,0 & \code{Yedge} & $y$方向の辺 \\
0,0,1 & \code{Zedge} & $z$方向の辺 \\
0,1,1 & \code{Xface} & 法線が$x$方向の面 \\
1,0,1 & \code{Yface} & 法線が$y$方向の面 \\
1,1,0 & \code{Zface} & 法線が$z$方向の面 \\
1,1,1 & \code{Cell}  & セル中心 \\
\end{yzntable}

位置を表で持たずにビットで符号化するのは、次の性質があるからである。\textbf{ある軸に沿って半セル移動すると、その軸のビットだけが反転し、ほかは変わらない}。コード中で配置型が変わる箇所は、すべてこの操作一つで表せる。添字の移動も同じビットから一行で求まる。$+\tfrac12$セル移動するときは、ビットが既に\code{1}なら添字を一つ増やす。$-\tfrac12$セル移動するときは、ビットが\code{0}なら添字を一つ減らす。手書きのオフセット表はないので、それをほかと整合させる手間もない。

1次元で意味をもつのは$x$のビットだけであり、2次元では$x$と$y$のビットである。未知数は1次元では\code{Xedge}に、2次元では\code{Zface}に置く。\val{"CZ"}モデル（\cref{sec:conc:cz}）の濃度は\code{Cell}に置く。いずれも面内のすべての軸について半セルずれた位置であり、\cref{ch:governing}でセル中心と呼んだ位置にあたる。

場は\code{Field<GRID,S,T>}として宣言する。配置型\code{S}は実行時のデータではなくテンプレート引数である。そのため配置を取り違えると、誤った答えが出るのではなくコンパイルエラーになる。辺中心の場を渡すべきところにセル中心の場は渡せない。\code{C\_FiniteDifferences::GradX}などに配置型ごとのオーバーロードがあるのも同じ理由による。場の微分は、どちら向きに半セル移動したかを型から知ることができる。

\section{ゴースト層}
\label{sec:disc:ghost}

格子は境界の外側に\textbf{各側二層のゴースト層}をもつ（\code{C\_Grid*::NumGhostX()}）。膜厚の方程式は四階であり、境界点で評価すると二つ隣の点まで参照するので、一層では足りない。

\key{Grid.Nx}は内部点だけを数える。したがって\key{Nx = 300}と宣言した格子は、$x$方向に$\mathrm{Nx} + 4 = 304$点をもつ。メモリを見積もるときや、デバッグ出力から行数を読み取るときは、この点に気をつける。

\code{CheckGhost(s)}は、すべての点を次の三つに分類する。

\begin{description}
  \item[\code{Inside}] 計算領域の内部。物理項が寄与を加えるのはここに限る。
  \item[\code{bXmin}, \code{bXmax}, \dots] 境界そのもの。値は境界条件で決まる。2次元では角も、3次元ではさらに辺もここに含まれる。
  \item[\code{Outside}] 最も外側のゴースト層。
\end{description}

\section{点から行列の行へ}
\label{sec:disc:rows}

行数は内部点の数ではなく、\textbf{ゴーストを含めた点の総数}$\mathrm{imax()} - \mathrm{imin()} + 1$である。点から行番号への写像は、次元ごとに\file{GridTraits.h}の\code{Index()}一箇所にまとめてある。写像を\cref{tab:row-index}に示す。

\begin{yzntable}{点から行への写像（次元ごとに一つ）\label{tab:row-index}}{@{}l l@{}}
\textsf{\bfseries 次元} & \textsf{\bfseries 行} \\ \midrule
1次元 & $\mathrm{row} = i_x$ \\
2次元 & $\mathrm{row} = n_x\,i_y + i_x$ \\
3次元 & $\mathrm{row} = n_x n_y\,i_z + n_x\,i_y + i_x$ \\
\end{yzntable}

ゴーストの行には、\code{CRS\_Matrix::SetAsBC()}で境界条件の行という印を付ける。印の付いた行は時間積分の対象外であり、\code{MakeMatrix\_IA()}と\code{MakeRHS()}はこれを飛ばす。物理項も\code{IsBC()}を見て、その行には加算しない。境界条件の行の係数は初期化時に一度だけ書き込む。その後行列をゼロクリアしても係数は残るので、作り直す必要はない。

ステンシルは1次元では五点であり、\code{M m o p P}と書く。中心の点とその両側の二点ずつである。2次元では十三点であり、九点の十字形に四つの対角点\code{mm}、\code{pm}、\code{mp}、\code{pp}を加える。\textbf{対角点は省略できない}。重調和作用素は交差項$\partial^4/\partial x^2\partial y^2$を含み、その差分係数は$1/(\Delta x^{2}\Delta y^{2})$のオーダーで、ステンシルはちょうどこの四つの対角点になるからである。対角点を除くと、$\nabla^{2}\nabla^{2}$の離散化ではなくなる。

各行の列を挿入する順序は、\code{T::FullIdx(k)}などの並びで決まっている。この順序は見た目の問題ではない。行列・ベクトル積での浮動小数点加算の順序がこれで決まるので、変えるとビット単位の再現性が失われる。\cref{ch:reproducibility}を参照されたい。

\section{セル界面の移動度と、その形を自由に選べない理由}
\label{sec:disc:positivity}

膜厚の方程式は四階の\emph{縮退}拡散方程式であり、移動度$D(h) = h^{3}$は液膜がなくなるところで消える。接触線を扱えるのはこの縮退があるからだが、不用意に離散化すると、同じ縮退のために膜厚が負になる。二つのセルの間のセル界面には、両側の膜厚から作った移動度が一つ必要である。その作り方によって、スキームが負の膜厚を生みうるかどうかが決まる。

鉛直一様モデルでは、Diez and Kondic (2002)の正値性保存補間（同論文の式(23)）を用いる。
%
\begin{equation}
  D_{k}^{(x)} \;=\; \frac{2\,h_{k}^{2}\,h_{k+1}^{2}}{h_{k} + h_{k+1}},
  \qquad
  D_{k}^{(y)} \;=\; \frac{2\,h_{k}^{2}\,h_{k+n_x}^{2}}{h_{k} + h_{k+n_x}} .
  \label{eq:disc:pps}
\end{equation}
%
これは場当たり的な平均ではない。一般的な処方
%
\begin{equation}
  D_{i} \;=\; \frac{h_{i+1} - h_{i}}{g_{i+1} - g_{i}},
  \qquad
  g(h) \;=\; \int \frac{\mathrm{d}h}{D(h)} ,
  \label{eq:disc:pps-general}
\end{equation}
%
を$D(h) = h^{s}$、$s = 3$として評価したものである。実装は\file{model/\bk well\_mixed/\bk LaplacePressureKernel.h}の\code{NonlinearTerm()}にある。

この形を選ぶ理由は次のとおりである。\cref{eq:disc:pps}は積$h_k h_{k+1}$と調和平均$2h_k h_{k+1}/(h_k + h_{k+1})$の積に分解できる。そのため片側が薄くなると$O(h^{2})$で消え、片側がゼロなら厳密にゼロになる。したがって乾いた領域の縁では、セル界面の流束がおのずと止まる。すぐに思いつく別案は算術平均$(h_k^{3} + h_{k+1}^{3})/2$である。しかしこれでは、完全に乾いた隣接セルに対しても$h_k^{3}/2$の移動度が残る。その結果、乾いた領域から液体が引き出され、$h$が負になる。

\begin{warn}[この式を単純化してはならない]
\cref{eq:disc:pps}を算術平均に書き換えると、正値性の保証が失われる。薄い側とともに消えないほかの式に書き換えても同じである。症状はクラッシュではなく、接触線近傍での負の膜厚として現れ、それが下流で粘度と濃度方程式を汚染する。この性質は\file{yuzen/tests/test\_positivity\_mobility.cpp}で固定している。このテストでは、\cref{eq:disc:pps}が\cref{eq:disc:pps-general}と一致すること、移動度が$O(h^{2})$で減少してゼロに達すること、算術平均がその両方を満たさないことを確認する。
\end{warn}

\subsection{ポテンシャルの各項に共通のセル界面移動度}
\label{sec:disc:positivity-one}

分離圧と鉛直重力には、同じセル界面移動度を用いる。ビット単位で\cref{eq:disc:pps}と同じ値である（\code{Mobility::Face}。どちらかの側が乾いている場合や、\val{"Pinning"}の辺のゴースト$-h$である場合はゼロ）。このときセル界面の流束は、一つの移動度とポテンシャル全体$\varphi = \hat\sigma\nabla^2(H+S) + \Dpn\,\Pihat(H) - \Bo\,(H+S)$の差との積になる（\cref{sec:output:potential}）。そのためスキームの散逸$\sum_f D_f(\Delta\varphi)^2/\Delta x$は二乗和になり、離散自由エネルギーは増加しない（\cref{sec:output:energy}）。定常な液膜は$\varphi$が一様な状態であり、離散エネルギーの最小点にあたる。起伏上の鉛直重力の項$D_f\,\Delta s$は、$K = D_f/(h_a + h_b)$として$K(h_a + h_b)\,\Delta s$の形で$A$に入る。セル界面の両側のセルの分担は等しい。以前はこの二つの項に算術平均$(h_a^3 + h_b^3)/2$を用いていたが、定常に近い液膜で\file{Energy.dat}に自由エネルギーの増加が現れたため改めた。この平均では、乾いた隣接セルに対して$h^3/2$の分離圧流束も残っていた。

\subsection{ナビエすべりがある場合のセル界面移動度}
\label{sec:disc:positivity-slip}

\key{Physics.Slip}（\cref{sec:ref-physics:slip}）を用いると、移動度は$D(h) = h^{3} + a h^{2}$、$a = 3\Lambda$となる。このとき\cref{eq:disc:pps-general}に用いる関数は次のとおりである。
%
\begin{equation}
  g(h) = -\frac{1}{a h} + \frac{1}{a^{2}}\ln\frac{h+a}{h}
       = -\frac{x - \ln(1+x)}{a^{2}}, \qquad x = \frac{a}{h} .
  \label{eq:disc:pps-slip}
\end{equation}
%
どちらの形もどこかで桁落ちを起こすので、\file{model/Mobility.h}はその範囲で級数に切り替える。$x = 0.1$未満では、$x - \ln(1+x)$を$\sum_{n=2}^{18}(-1)^n x^n/n$として和をとる（剰余は相対値で$5\times10^{-19}$未満）。二つの高さの差が平均$\bar h$の$10^{-3}$未満の場合は、セル界面移動度を$D(\bar h)/(1 + \delta^{2}K/24)$とする。ここで$\delta = h_{k+1} - h_k$、$K = D\,(1/D)''$であり、打ち切り誤差は$O((\delta/\bar h)^4) < 10^{-12}$である。桁落ちのない参照値と比べた最大相対誤差は$4\times10^{-13}$であり、各切り替えの前後で値は$10^{-12}$まで連続である。すべりがなければこれらの処理には入らず、\cref{eq:disc:pps}を式のとおりに評価する。分離圧と鉛直重力には、すべりの有無によらず同じセル界面移動度を用いる（\cref{sec:disc:positivity-one}）。

セル界面移動度は乾いた側で消えることに変わりはない。ただしその近傍では$D \sim a h^{2}$なので、消え方は$O(h^{2})$ではなく$O(h)$になる。

\begin{warn}[すべりは正値性の保証を弱める]
すべりがあると、液膜を正に保つ縮退が一次弱くなる。液膜が分離圧で支えられている場合や、高さの一定割合（およそ$0.1\,h_0$）より薄くならない場合には、目に見える影響はない。一方、分離圧なしで膜厚ゼロに向かう液膜は、すべりがなければ負にならない場合でも負になりうる。\file{History.dat}の\code{HHmin}を監視する必要がある。
\end{warn}

\subsection{\texorpdfstring{\val{"CZ"}}{CZ}モデルでも同じ移動度を用いる}
\label{sec:disc:positivity-concz}

\val{"CZ"}モデルは濃度、したがって粘度の厚さ方向の分布を解くので、面内速度は解析的に求まらない。\file{model/ConcZ/CZ2D/CZ\_HeightEquation\_2D.cpp}の\code{CalVelocity()}で$1/\mu$と$z/\mu$を液膜の厚さ方向にシンプソン則で積分し、濃度方程式はその速度で移流する。

一方、膜厚そのものは鉛直一様モデルと同じ保存形
%
\begin{equation}
  \frac{\partial h}{\partial t} = -\frac{\partial Q}{\partial x} - E ,
  \label{eq:disc:cz-conservative}
\end{equation}
%
で進める。セル界面の流束には同じ移動度\cref{eq:disc:pps}を用いるが、深さ積分した流束は$h^{3}$ではなく$h^{3}/3\mu$なので、三で割る。粘度が一様な極限では、\val{"CZ"}の速度分布を厚さ方向に積分するとこの係数が厳密に得られる。実装は\code{FaceMobility()}である。鉛直一様モデルで同じ$1/3$にあたるのは、毛管項の\code{\_Magnitude}である。

\begin{note}[セル界面での解像された粘度]
セル界面を通る流束は、解像された速度分布の深さ積分
$Q = h^3 b\int_0^1 (1-\zeta)^2/\mu\,d\zeta + h^2\tau\int_0^1 (1-\zeta)/\mu\,d\zeta$
である。ここで$b$は駆動する圧力勾配と体積力、$\tau$はマランゴニ応力である。したがって移動度は$1/\mu(c(\zeta))$を$3(1-\zeta)^2$で重み付けした平均を含み、マランゴニ流束は$2(1-\zeta)$で重み付けした平均を含む。表面張力には表面濃度での値を用いる（\code{FaceInvViscosity()}、\code{FaceSurfConc()}）。正値性を保つ$h$の因子は変わらず、粘度が一様なら厳密に$h^3/3\mu$になる。セル界面ごとに深さ平均濃度での粘度と表面張力を一つずつとる移動度では、液膜を動かす流束が、溶質を運ぶ速度分布の積分と一致しない。濃度が成層している場合、この誤差は格子を細かくしても減らない（\cref{sec:verification:mms}）。
\end{note}

両モデルはどちらも保存形で、移動度も共通なので、濃度が厚さ方向に一様な極限では一致する。両者は異なる経路でこの極限に至るため、この一致は意味のある検査になる。比較は\file{tools/concz\_vs\_wm.py}で行い、テストスイートには\code{concz\_vs\_wm}として含めている。\key{Fluid.Viscosity.Exponent}を\val{0}とし、蒸発と溶質マランゴニ効果を無効にすると、濃度は$h$に影響しなくなる。このとき二つのモデルの最大膜厚は相対差$10^{-6}$以内で一致する。体積はどちらも厳密に保存され、二つの\file{History.dat}で同じ値になる。

\begin{note}[この比較での二つの落とし穴]
第一に、適応時間刻みを有効にしたままだと、二つのモデルは異なる時間刻みの履歴をたどり、出力時刻が次第にずれる。この状態で各\file{History.dat}の最終行を比べると、異なる二つの時刻を比べることになる。生じる数パーセントの差は格子細分化に対して非単調に変動し、離散化誤差と見分けがつかない。時間刻みは固定する必要がある。第二に、出力間隔は\key{EndTime}より十分小さくする。そうしないと共通の時刻が初期条件だけになり、両モデルが完全に一致しているように見えてしまう。
\end{note}

\section{対流スキーム}
\label{sec:disc:convection}

対流的な項は三種類ある。濃度の移流と、鉛直一様モデルの膜厚の方程式に現れる面内体積力およびマランゴニ流束である。後の二つは、移動度と勾配の積ではなく、速度で運ばれる膜厚の形で入る。いずれもセル界面での場の値が必要であり、セル界面の周りの四つのセル$f_{LL}, f_L, f_R, f_{RR}$から求める。固定スキームでは、この値を
%
\begin{equation}
  f_{\text{face}} = \sum_j c_{1,j} f_j + \operatorname{sign}(u)\sum_j c_{2,j} f_j ,
  \label{eq:disc:face}
\end{equation}
%
とする。UTOPIAでは$c_1 = (-1, 7, 7, -1)/12$、$c_2 = (-1, 3, -3, 1)/12$であり、河村・桑原スキームでは$c_2$がその三倍になる。どちらも三次精度であり、移流の主要誤差はそれぞれ$-|u|\Delta x^3/12\,\partial^4 f$、$-|u|\Delta x^3/4\,\partial^4 f$である。\val{"Central4"}には$c_2$がない。セル界面値は四次精度だが、発散は二次精度にとどまり、散逸はない。UTOPIAはLeonard, MacVean and Lock (1993)による一様三次精度の多項式補間である。河村・桑原スキームはKawamura, Takami and Kuwahara (1986)のスキームをセル界面値の形に書き直したもの、\val{"QUICK"}はLeonard (1979)の二次風上補間である。いずれも参考文献に挙げた。

\val{"Symmetry"}または\val{"Pinning"}の辺のセル界面では、膜厚の方程式で風上化した三つの項は流束をもたない。そこではこれらの項の移動度を、$A$、ヤコビアン、速度のいずれでもゼロとする（\code{UpwindKernelBase::WallFace}）。これらの流束は勾配の形ではないので、鏡像のゴーストを置いてもセル界面値はゼロにならないからである。

リミッター付きスキームでは、セル界面値を風上側から書く。$+x$方向の流れでは$(UU, U, D)$をセル$(LL, L, R)$、$-x$方向の流れでは$(RR, R, L)$とする。二つの差を$a = f_U - f_{UU}$、$b = f_D - f_U$とすると、
%
\begin{equation}
  f_{\text{face}} = f_U + \tfrac12\,\Phi(a, b), \qquad
  \Phi_{\text{Koren}} = \operatorname{sign}(b)\min\!\left(2|a|,\ \tfrac{2|b| + |a|}{3},\ 2|b|\right), \qquad
  \Phi_{\text{vL}} = \frac{2ab}{a + b},
  \label{eq:disc:limited}
\end{equation}
%
であり、$a$と$b$が同符号でなければ$\Phi = 0$とする。これらはKoren (1993)とvan Leer (1974)のリミッターである。ゼロになりうる勾配で割ることはない。Korenのリミッターは、$0.4 \le a/b \le 4$の範囲ではUTOPIAと同じであり、それ以外では有界である。セル界面値は$f_U$と$f_D$の間から出ず、極値では$f_U$に戻る。有界性と引き換えに、滑らかな極値は局所的に一次精度に削られ、滑らかな解の精度が落ちる。正弦波では、最大誤差の収束次数は3ではなく約1.5になる。ステンシルを鏡映しても値は最後のビットまで同じなので、\cref{sec:verification:equivalence}の検査はすべてのスキームで成り立つ。

濃度方程式は、係数を固定すれば$c$について線形である。この方程式では、リミッター付きのセル界面値を最新の反復値で凍結する。セル界面の風上セルの行では$\chi = \Phi/a$として$f_U + \chi/2\,(f_U - f_{UU})$と書き、それ以外の行では$\psi = \Phi/b$として$f_U + \psi/2\,(f_D - f_U)$と書く。こうすると凍結した行列はM行列になり、クランク・ニコルソン法の陽的な半分はクーラン数1まで正になる。そのためクーラン数1までの範囲では、ステップを進めても$c$は範囲内に保たれる。ただし時間刻みをこの範囲に制限するわけではない。範囲を超えたステップは、この方程式でも膜厚の方程式でも\file{Warnings.txt}に報告する（\cref{sec:ref-numerics:convection}）。鉛直一様モデルの濃度はニュートン法のパスごとに解き直すので、凍結した値は液膜とともに収束する。ステップの完了時には、リミッターをその解自身で凍結する（\cref{sec:disc:wm-conservative}）。各ステップの陽的な半分は、ステップ開始時の状態と、その状態でのリミッターから作る。\val{"CZ"}では、各段はその段で評価する濃度を用いる。これはリミッター付き作用素の中点則にあたる。

\paragraph{溶質を切り捨ても補いもせずに$0 \le c \le 1$とする。}濃度は質量分率であり、完了した各ステップで$0 \le c \le 1$の範囲に収める。蒸発を溶媒に限る場合（\cref{sec:terms:evaporation}）、1を超える原因は丸め誤差と、急峻な前線で非有界スキームが生じるオーバーシュートとアンダーシュートだけである。KK、UTOPIA、Central4、QUICKではどのクーラン数でも、Korenとvan Leerではクーラン数1を超えるとこれが起こる。鉛直一様モデルでは、実セルが1を超えて保持する分$h(c-1)$を、セル界面を共有する隣接セルに、それぞれの余裕$h(1-c)$に比例して渡す。受け取りきれなかった分は、それが置かれたセルから次のパスでさらに渡す（最大四回）。それでも残った分は、余裕のあるすべてのセルに余裕に比例して配分する。0を下回るセルは、逆向きに同じ方法で補う。不足分$-hc$を、セル界面を共有する隣接セルから、各セルが保持する溶質$hc$に比例して、最大四回のパスで受け取る。その後は溶質を保持するすべてのセルから受け取る。セルの面積は等しいので、\code{ResinVolErr}が収支をとる$hc$の和（\cref{sec:output:history}）は丸め誤差の範囲で保たれる。ゴーストは移動後のセルから改めて作る。各パスでは何かを動かす前に全セルの余裕を求めるので、結果はセルの順序によらない。$[0, 1]$を出るセルがなければ何も変更しない。残余が生じるのは、全体として液体より多くの溶質を保持する液膜か、溶質がゼロ未満の液膜だけである。その場合は切り捨てるか補い、\file{Warnings.txt}に報告する。実行全体で移した溶質の量は\file{JobInfo.txt}に示す（\cref{sec:output:jobinfo}）。0未満のまま放置した場合、$c = 0.7$から$0.005$への前線のアンダーシュートは$-0.04$から$-0.09$に達した。QUICKとCentral4、およびクーラン数1.5のUTOPIAでは、発散判定の$-0.1$を越えて実行が停止した（発見済み、修正済み）。ニュートン法のパスの中では、粘度、マランゴニ応力、蒸発の計算に用いる作業用の場を1で頭打ちにする。一方、Korenとvan Leerのリミッターは、方程式を解いたままの反復値で凍結する。頭打ちにした値で凍結すると、一つのセル界面を共有する二つの行が異なるセル界面値をとることになり、そこで移流が流束形式でなくなるからである。\val{"CZ"}では、セルの超過分を柱の中で一つ下のセルへ渡す。柱のセルはどれも同じ体積をもつ。柱の平均が1を超える場合は、警告を出して$c$を1より大きいまま保つ（溶媒に限る蒸発のため、実際には起こらない）。0を下回るセルは、同じ柱の中で溶質を保持するセルから、その溶質に比例して補う。平均が0を下回る柱は、液膜に沿った移流のアンダーシュートで生じる。この場合は不足分を隣接する柱から、それぞれの溶質に比例して受け取る（柱全体をまとめてスケーリングするので、隣接する柱のどのセルも0を下回らない）。これを最大四回のパスで行い、その後は鉛直一様モデルのセルと同様に、溶質を保持するすべての柱から受け取る（\code{CZ\_HeightEquation\_2D::RedistributeDeficit}）。0未満のまま放置した場合、クーラン数0.4のUTOPIAで計算した\val{"CZ"}の0.7から0.005への前線は、発散判定を越えて実行が停止した（発見済み、修正済み）。

膜厚の方程式では、行列$A$に現在の液膜で凍結したセル界面の重みを保持し、$A(h)h$が厳密に流束となるようにする。ヤコビアンには、リミッター付きのセル界面膜厚$h_f$の微分の残りの部分を加える。$g = \partial h_f/\partial h$として、体積力では$w(3g - w_t)h_f^2/\mu$、マランゴニ流束では$(2g - w_t)h_f\,\Delta c/\mu$である。これによりニュートン法は二次収束を保つ（\cref{sec:verification:schemes}）。Korenの$\Phi$の折れ点では片側微分になる。マランゴニ流束は、$\Delta c$だけの符号ではなく、流束自身の向き、すなわち$\sigma_c\,\Delta c$の符号で風上化する。$\Delta c$だけで判定すると、\key{Fluid.SurfaceTension.Coef}が負のときに風下側を選んでしまう。

\subsection{鉛直一様モデルの濃度の保存形}
\label{sec:disc:wm-conservative}

鉛直一様モデルの濃度方程式は、各セルの溶質$hc$のステップ間の有限体積収支として解く。
%
\begin{equation}\label{eq:disc:wm-conservative}
  \begin{aligned}
    h^{1}c^{1} - h^{0}c^{0}
      = {}& -\Delta t \sum_f s_f \bigl[\theta\,G_f^{1} + (1-\theta)\,G_f^{0}\bigr]
        + \Delta t\,S\,c_S \\
      &{} - \Delta t\,\bigl[\theta\,r^{1}c^{1} + (1-\theta)\,r^{0}c^{0}\bigr],
    \qquad
    G_f = F_f\,c_f - D(c_f)\,h_f\,\frac{c_{\mathrm{hi}} - c_{\mathrm{lo}}}{\Delta x},
  \end{aligned}
\end{equation}
%
ここで$s_f = \pm 1/\Delta x$は、セル界面$f$の外向きの符号をそれを挟むセルの間隔で割ったものであり、$\theta = 1/2$である。$F_f$は膜厚の方程式がそのセル界面に通す体積流束である。有効な各項がそのセル界面について$A$に保持する値の和であり、毛管、鉛直重力、分離圧の流束と、風上化したセル界面膜厚を用いる面内体積力、遠心力、マランゴニ応力の流束からなる（各カーネルの\code{FaceFlux}）。$c_f$は\key{Numerics.Convection.Concentration}によるセル界面値（\cref{eq:disc:face}）であり、$F_f$の符号で風上化する。$D(c_f)\,h_f$には二つのセルの平均を用いる。$S$はノズルから供給される液体であり、膜厚の方程式と同じくステップの中間で評価する。$c_S$は供給された液体が到着時にとる濃度である。\key{Physics.Source.ConcType = "local"}（既定値）では$c_S = \theta c^{1} + (1-\theta)c^{0}$、すなわち着地点の濃度とする。\val{"supplied"}では$c_S = \key{Physics.Source.Conc}$とし、新しい液体をセルの柱の厚さ全体に混合する。流れのない液膜では$c^{1} = (h^{0}c^{0} + \Delta t\,S\,c_S)/(h^{0} + \Delta t\,S)$となる。$r$は1次元の遠心項による側方への流出であり、そのセルの濃度で出ていく。蒸発は溶媒だけを奪うので式には現れず、$h^{1}$に含まれる。この式は\cref{eq:conc-dim}そのものである。$c$についての\cref{eq:conc-adv}の形も連続系では一致するが、離散系で溶質を保存するのは収支の形だけである。\cref{eq:conc-adv}の形で$c$について解いた場合、閉じた周期液膜上の0.9から0.3への前線を流れで運びつつラウール則で乾燥させると、領域を0.3回横断する間に溶質の$1\times10^{-5}$から$3\times10^{-4}$が失われた。損失量はスキームと格子（50から200セル）によって異なる（発見済み、修正済み。現在は1次元で$4\times10^{-16}$から$3\times10^{-15}$、2次元で$2\times10^{-14}$。\cref{sec:verification:analytic}の\code{conc\_drying}）。

方程式はステップでの変化量$d = c^{1} - c^{0}$について解く。液膜自身のステップ$h^{1} - h^{0} = -\Delta t \sum_f s_f[\theta F_f^{1} + (1-\theta)F_f^{0}]
- \Delta t\,[\theta r^{1} + (1-\theta)r^{0}] - \Delta t\,(E - S)$に$c^{0}$を掛けて\cref{eq:disc:wm-conservative}から差し引くと、
%
\begin{equation}\label{eq:disc:wm-change}
  \begin{aligned}
    &h^{1} d + \theta\Delta t \sum_f s_f\bigl[F_f^{1} d_f - K_f^{1}(d_{\mathrm{hi}} - d_{\mathrm{lo}})\bigr]
       + \theta\Delta t\,(r^{1} - S)\,d \\
    &\quad= -\Delta t \sum_f s_f \bigl[\theta F_f^{1} + (1-\theta)F_f^{0}\bigr](c^{0}_f - c^{0})
       + \Delta t\,E\,c^{0} \\
    &\qquad\quad{}+ \Delta t \sum_f s_f \bigl[\theta K_f^{1} + (1-\theta)K_f^{0}\bigr](c^{0}_{\mathrm{hi}} - c^{0}_{\mathrm{lo}}),
  \end{aligned}
\end{equation}
%
が残る。ここで$K_f = D(c_f)h_f/\Delta x$である。各セル界面値$c^{0}_f - c^{0}$は、セル界面の重みと$c^{0}_k - c^{0}$の積の和として作る。重みの和は1なので、$c^{0}$が一様で蒸発がなければ右辺は厳密にゼロになり、$d = 0$となる。蒸発の項$\Delta t\,E\,c^{0}$は、ほかに何も作用しなければ$h^{1}c^{1} = h^{0}c^{0}$となるように溶質を濃縮する。\key{EnableConvection}を無効にすると$c_f - c$の項は落ちるが、それでも溶質のある液膜は薄くなったり厚くなったりする（$\theta\Delta t \sum_f s_f
F_f^{1} d$）。このとき流れは液体だけを動かし、中の溶質は動かさないので、溶質は収支の形にならない。ステップ終了時の値は$c^{0} + d$である。

液膜$h^{1}$は、同じステップのニュートン反復から得る。各パスでは、そのパスの液膜とその流束を用いて濃度を解く（\cref{sec:time:coupling}）。ただし、これらの流束がその液膜と整合するのは反復が収束した後に限られる。反復の後でステップを完了する（\file{HeightEquationBase.h}の\code{AfterNKLoop}）。液膜には反復が収束した値を用いる。流束、側方流出、供給はその液膜について最後のパスの濃度で評価し、蒸発はそのパスで用いた値とする。これらの発散では、液膜が実際に進んだステップにセルごとにまだ足りない分があり、これがその液膜でのステップの残差である。残差のセル平均はステップの体積誤差であり、そのまま残す。平均を差し引いた残りは、内部のセル界面の流束に累積和として加える。まず$y$方向に、各セル行の和をその行に均等に配分し、次に各行の中で残りを$x$方向に配分する。こうするとセル界面流束は、液膜が進んだステップと厳密に整合する（\code{ConsistentFluxes}）。辺のセル界面や壁には手を触れず、何かを解くこともない。仮に流束から液膜を求めると、すなわち$h^{0}$から流束の発散を引くと、収束した液膜とは上の残差だけ異なる。セルよりはるかに厚い液膜では、この残差は巨大な流束の丸め誤差の水準にある（$\epsmach\,\Delta t/\Delta t_{\mathrm{explicit}}$、\cref{sec:convergence:floor}）。これを液膜として採用すると、\file{nozzle\_1d\_rotating}では膜厚が負になる。続いて、これらの流束とこの液膜で濃度をもう一度解き、各セルの溶質を\cref{eq:disc:wm-conservative}から作る。各セル界面の溶質流束は、どちら側のセルが求めても、その周りの四つのセルから一度だけ作る（\file{ConcEquationBase.h}の\code{Reconstruct}）。これにより、全セルにわたる$hc$の和は、線形求解の許容誤差によらず、辺を通る流束と湧き出しの分だけ、丸め誤差の範囲で変化する。\file{History.dat}が\val{"FreshInlet"}や\val{"Outlet"}を横切る量として数えるのは、この辺の流束である。すなわち各辺のセル界面の体積流束（\code{ConsistentFluxes}はこれに手を触れない）と、\code{Reconstruct}がそこで作った溶質流束である（\cref{sec:output:history}）。リミッター付きスキームでは、リミッターをその解自身で凍結して、値が落ち着くまで（$10^{-12}$、最大八回）解き直す。これにより、ステップのセル界面値は解のリミッター付きセル界面値と一致する。その後$c$を$[0, 1]$に収める（\cref{sec:disc:convection}）。パスまたはステップの終わりで液膜に液体のないセル（$h \le 0$。液膜に対して時間刻みが長すぎる場合に生じる）があるときは、濃度を前のパスの値のままにする。そこでは$hc$の収支が成り立たないので、ステップは液膜の側で失敗させる（発見済み：以前は濃度の求解が先に失敗し、濃度ソルバーの名を挙げて実行が停止していた。修正済み）。そのようなステップでは、液膜のニュートン修正量が収束から遠いところで停滞する。その結果ステップが縮められるか、固定刻みでは停滞を報告して実行が停止する（\file{tools/input\_guards.py}の\code{wm\_diverges}、\cref{sec:convergence:stagnation}）。それでも負の膜厚、たとえばファイルから読み込んだ負の液膜は、発散として実行を停止させる（\code{wm\_film\_below}）。

強い面内流れ、ノズル、鉛直重力、毛管圧、さらに1次元では回転があっても、一様な濃度は$6\times10^{-16}$まで一様に保たれる（\code{conc\_uniform}）。一般には、液膜に対する溶質流束の丸め誤差$\epsmach\,\Delta t\sum_f|F_f|/h$の水準まで保たれる。この誤差は、流束から改めて作った溶質に持ち越される（\file{tools/proptest.py}では、\val{"Pinning"}の辺に隣接する薄い液膜で最大$3\times10^{-12}$）。モデルのすべての項を有効にすると、製造解は$\Delta x$と$\Delta t$について二次で収束し（\cref{sec:verification:mms}）、\cref{sec:verification:schemes}の斜行波は1.91と1.97で収束する。

\subsection{\texorpdfstring{\val{"CZ"}}{CZ}の濃度の保存形}
\label{sec:disc:cz-conservative}

\val{"CZ"}の濃度方程式は、$\zeta$格子上の$hc$の有限体積収支である。膜厚の方程式は各セル界面を通る流束を保持し、ステップの終わりの膜厚をその発散で表す。濃度方程式は、各流束を層ごとの流束に分ける。層ごとの流束は解像された速度分布のセル平均であり、柱の各層の和がセル界面ごとに膜厚の流束と一致するように調整する。$\zeta$方向のセル界面を通る流束は、格子の離散的な保存則から、基板側から順に積み上げて求める。蒸発する溶媒は表面のセル界面から出ていく。表面の上と基板の下のゴーストセルは、そこでの全流束（移流と、斜面に沿う成分を含む拡散流束全体）がゼロになるように設定する。ステップは$c$の変化量について解くので、一様な$c$は厳密に変化しない。線形ソルバーが残す小さな溶質の残差は、変化量の一様なシフトとして戻す。シフトの大きさは、シフト自身が表面項を通して送り出す溶媒の分を見込んで決め、ゴーストにも同じシフトを適用する（周期の継ぎ目では、両側がそれぞれ自分の側からそこでの流束を作る）。これによりステップの溶質は丸め誤差の範囲で収支が合う。リミッター付きスキームでは凍結した風上形と風下形を用いるので、$c$は有界に保たれる。リミッターがステップの中点で落ち着くまで解き直す。

\key{Numerics.Convection.Concentration}のスキームは、$x$、$y$、$\zeta$のいずれのセル界面にも同様に用いる。中心スキームはどの方向にも散逸をもたないので、格子が解像できる最短の波長では、解はモデル本来の性質をそのまま示す。$\sigma_c\,\partial c/\partial\zeta
> 0$で成層し拡散係数の小さい液膜は、その波長で不安定である。中心スキームではこれらの波が成長するが、風上寄りのスキームでは格子が十分粗ければ減衰する（\cref{sec:verification:schemes}）。
